(a) The gradient is $(y^2,2xy,-2z)$, so the singular locus is the line $y=z=0$. At $(a,0,0)$ with $a\ne0$, the quadratic term is $ay^2-z^2$, a product of two distinct linear forms. At the origin it is the double plane $z^2=0$. Over $\mathbb{R}$, the sections $x=a>0$ are the two lines $z=\pm\sqrt a\,y$, which meet at the pinch point as $a\to0$.

(b) The gradient of $x^2+y^2-z^2$ is $(2x,2y,-2z)$, so only the origin is singular. The surface is the quadratic cone, with a conical double point at its vertex.

(c) The singularity equations are $y+3x^2=x+3y^2=0$. Multiplying by $x,y$ and using $xy+x^3+y^3=0$ gives $x^3=y^3=0$. Thus the singular locus is $x=y=0$. Since the equation is independent of $z$ and its quadratic term is $xy$, this is a double line, locally the product of a node with a line.
